\chapter{Podsumowanie i~wnioski}

Celem pracy było opracowanie i~zaimplementowanie w~języku opisu sprzętu wybranych etapów toru nadajnika PDSCH w~standardzie 5G NR, ze szczególnym uwzględnieniem dopisania sum kontrolnych CRC oraz kodowania QC-LDPC, przy założeniu stałej magistrali danych 64~bity. Zakres obejmował również przygotowanie danych do kodowania, minimalny podzbiór operacji powiązanych z~\english{rate matching'em} oraz weryfikację bitowo-dokładną w~oparciu o~wektory referencyjne. 

\section*{Ocena realizacji celów}
Na podstawie przeprowadzonych prac można stwierdzić, że założone cele zostały osiągnięte w~zakresie przyjętym w~pracy:

\begin{itemize}
  \item Zaimplementowano kompletny łańcuch przetwarzania zintegrowany w~module top-level \texttt{PDSCHTx}, realizujący kolejne etapy: dopisanie CRC do TB, przygotowanie danych do kodowania LDPC oraz kodowanie LDPC, z~emisją danych wyjściowych jako strumień 64-bit (CB po CB, a~znacznik końca TB pojawia się tylko raz na końcu całego strumienia). 

  \item Zachowano spójny format danych w~całym torze: interpretację bitów w~kolejności MSB-first oraz deterministyczny strumień 64-bit niezależnie od długości TB, z~jednoznaczną obsługą dopełnień poprzez metadane wyjściowe \texttt{out\_pad\_left} i~\texttt{out\_pad\_right}. 

  \item Opracowano środowisko weryfikacji oparte o~symulację HDL (Questa), testbench’e oraz referencję MATLAB, umożliwiające automatyczne porównanie wyników DUT z~wektorami wzorcowymi w~sposób bitowo-dokładny, wraz z~generacją plików wynikowych i~raportowaniem rozbieżności. 

  \item Zweryfikowano poprawność kluczowych mechanizmów wymaganych do zgodności z~przyjętym modelem referencyjnym: obliczenie oraz dołączenie CRC, przygotowanie oraz kodowanie LDPC, nieuwzględnianie \english{filler bits} w~strumieniu transmisyjnym oraz pominięcie dwóch pierwszych bloków długości $Z_c$ na wyjściu kodera dla każdego CB (mechanizm zgodny z~ideą \english{puncturing’u} w~procedurach 5G NR). 
\end{itemize}

\section*{Wnioski z~weryfikacji i~analizy implementacyjnej}
Wyniki rozdziału weryfikacyjnego pozwalają sformułować następujące wnioski praktyczne:

\begin{itemize}
  \item \textbf{Model pracy toru tworzy niepotrzebne opóźnienie:} w~aktualnej architekturze etap LDPC\_encode przechodzi do fazy emisji dopiero po przetworzeniu wszystkich bloków kodowych, co powoduje dominację tego modułu w~całkowitej latencji i~ma bezpośredni wpływ na opóźnienie end-to-end. 

  \item \textbf{Osiągalna przepustowość jest rzędu dziesiątek Mb/s:} na podstawie zmierzonej latencji (w cyklach) oraz oszacowanego $F_{\max}$ wyznaczono przepustowość goodput TB. Dla krótkich ramek narzut sterowania dominuje, natomiast dla dłuższych przypadków przepustowość stabilizuje się na poziomie kilkudziesięciu Mb/s. 
  \item \textbf{Projekt jest realizowalny w~docelowym FPGA z~istotnym zapasem zasobów:}
  po pełnej kompilacji w~Quartus Prime projekt mieści się w~układzie Cyclone~V z~wykorzystaniem ok. 29\% ALM oraz niskim udziałem pamięci i~bloków DSP, co pozostawia przestrzeń na dalszą rozbudowę. 

  \item \textbf{Ograniczeniem jest ścieżka krytyczna i~wymaganie 100 MHz:}
  analiza TimeQuest dla constraintu 100~MHz wykazała ujemny slack dla setup w~najgorszym narożniku, co przekłada się na $F_{\max}$ na poziomie ok. 70--71~MHz.
  Jednocześnie sprawdzenia hold oraz minimalnej szerokości impulsu są spełnione. 
\end{itemize}

\section*{Ograniczenia aktualnej wersji}
Należy podkreślić, że praca realizuje funkcjonalności do poziomu umożliwiającego bitowo-dokładną weryfikację i~ocenę realizowalności w~FPGA, jednak pełny tor PDSCH w~rozumieniu standardu obejmuje dodatkowe etapy. W szczególności:
\begin{itemize}
  \item nie zaimplementowano bloku mieszania (\textit{scrambling}), który był rozważany jako element opcjonalny,
  \item zrealizowano jedynie minimalny podzbiór operacji powiązanych z~\english{rete matching'em}, bez pełnej procedury dopasowania długości do zasobów transmisji,
  \item interfejs sterowania przepływem danych bazuje na \texttt{valid/last} bez mechanizmu handshake,
  \item system w~obecnym modelu obsługuje przetwarzanie jednego TB w~jednym cyklu pracy, a~rozpoczęcie kolejnego wymaga resetu układu,
  \item kodowanie LDPC jest realizowane w~sposób sekwencyjny, bez potokowania ani równoległego przetwarzania wielu bloków kodowych. 
\end{itemize}

\section*{Kierunki dalszych prac}
W oparciu o~uzyskane rezultaty oraz zidentyfikowane ograniczenia można wskazać następujące, logiczne kierunki rozwoju projektu:

\begin{itemize}
  \item \textbf{Rozszerzenie toru o~\english{scrambling} i~pełny \english{rate matching}} oraz ewentualnie o~dalsze elementy toru PDSCH. 

  \item \textbf{Usprawnienie architektury pod kątem latencji i~przepustowości:} Zrównoleglenie obliczeń, ograniczenie dominacji buforowania (np. poprzez częściowe potokowanie lub wcześniejszą emisję fragmentów strumienia) oraz redukcję logiki kombinacyjnej na ścieżkach sterujących w~celu podniesienia $F_{\max}$. 

  \item \textbf{Optymalizacja przekazywania danych na wyjście:}
  obecna organizacja emisji zakłada wyrównywanie fragmentów (m.in. granic CB oraz bloków parzystości) do słów 64-bit
  i~nierozpoczynanie kolejnych fragmentów w~środku słowa, co upraszcza sterowanie i~adresowanie buforów,
  jednak generuje puste bity w~słowach brzegowych oraz dodatkowe narzuty wynikające z~dopełnień.
  Naturalnym kierunkiem rozwoju jest zastosowanie bardziej zwartej emisji z~możliwością „sklejania” fragmentów
  bitowo (ciągły strumień bez sztucznych granic słowa), przy jednoczesnym zagęszczeniu buforów.
  Pozwoli to zmniejszyć objętość danych wyjściowych i~poprawić efektywną przepustowość,
  kosztem bardziej złożonej logiki.



  \item \textbf{Rozszerzenie interfejsu o~pełny \english{handshake}} (np. \texttt{AXI4}) i~tryb pracy ciągłej, umożliwiający obsługę kolejnych TB bez konieczności resetu. 

  \item \textbf{Rozbudowa pakietu testów:} zwiększenie pokrycia o~dodatkowe konfiguracje parametrów LDPC oraz testy regresyjne.
\end{itemize}

\section*{Podsumowanie końcowe}
Zrealizowany projekt potwierdza możliwość implementacji kluczowych elementów kodowania kanałowego dla toru PDSCH w~5G NR w~postaci modularnego systemu RTL o~stałej szerokości danych 64~bity, z~zachowaniem deterministycznego formatu strumienia i~jednoznacznej obsługi metadanych. Wyniki kompilacji w~FPGA wskazują na dużą rezerwę zasobów, natomiast analiza czasowa i~estymacja goodput pokazują, że aktualne ograniczenia wydajności wynikają głównie z~architektury modułu \texttt{LDPC\_encode} oraz czasu ścieżki krytycznej. Wskazane kierunki rozwoju stanowią bezpośrednią ścieżkę do rozszerzenia pracy w~stronę pełniejszej realizacji toru PDSCH oraz poprawy parametrów czasowych. 
